% chapters/16/section-03.tex
\section{Pemeriksaan Otomatis Tanpa Pemain}
\label{sec:uji-otomatis}

\subsection{Soft-lock: Kegagalan yang Tidak Terlihat}

Satu jenis kesalahan pada level tidak menimbulkan pesan apa pun, tidak
menghentikan program, dan tidak terlihat pada satu kali percobaan. Kesalahan
itu disebut \istilah{soft-lock}: keadaan ketika pemain tidak dapat melanjutkan
dan tidak dapat kembali, walaupun permainannya sendiri masih berjalan
\cite{unity-leveldesign}.

Pemain yang terjatuh ke dalam lubang yang terlalu tinggi untuk ditinggalkan
mengalami \istilah{soft-lock}. Programnya tidak berhenti, tidak ada kesalahan
yang tercetak pada jendela \istilah{Console}, dan penguji yang mencoba level itu
satu kali mungkin tidak pernah menemukan lubang tersebut. Yang mengalami
masalah hanya sebagian pemain, dan itulah yang membuat kesalahan ini mahal.

Karena kesalahan ini bersifat pasti, bukan soal selera, ia dapat diperiksa
secara pasti pula: dengan memeriksa setiap tempat pemain dapat berdiri, lalu
menanyakan apakah dari sana titik tujuan masih dapat dicapai.

\subsection{Jangkauan Lompatan Dihitung, Bukan Ditebak}

Pemeriksaan itu menuntut satu angka: seberapa jauh pemain dapat melompat.
Angka itu tidak perlu ditebak, sebab Anda sudah memiliki besarannya sejak
Bab~VI dan Bab~XIII. Tinggi lompatan tertinggi berasal dari hukum gerak:

\[
h_{\text{maks}} = \frac{v_0^2}{2g},
\qquad
t = \frac{v_0 + \sqrt{v_0^2 - 2g\,\Delta h}}{g},
\qquad
x_{\text{maks}} = v_{\text{jalan}} \cdot t .
\]

Dengan $v_0 = 700$ piksel per detik, $g = 1800$ piksel per detik kuadrat, dan
kecepatan jalan 300 piksel per detik, diperoleh tinggi 136,1 piksel dan jarak
mendatar 233,3 piksel. Pada peta yang petaknya berukuran 32 piksel, kedua
angka itu berarti:

\begin{table}[htbp]
\centering
\caption{Jangkauan lompatan yang dihitung dari besaran fisika Bab VI}
\label{tab:jangkauan-lompatan}
{\setlength{\tabcolsep}{4pt}\footnotesize%
\begin{tabularx}{\textwidth}{@{}L{7.0cm} Y@{}}
\toprule
\headrow \textbf{Ukuran} & \textbf{Nilai} \\
\midrule
Tinggi lompatan tertinggi & 136,1 piksel, yaitu 4,25 petak \\
\addlinespace
Jarak mendatar terjauh pada ketinggian yang sama & 233,3 piksel, yaitu 7,29 petak \\
\addlinespace
Jarak mendatar terjauh pada penurunan 5 petak & 288,6 piksel, yaitu 9,02 petak \\
\bottomrule
\end{tabularx}}
\end{table}

Jarak mendatar terjauh bergantung pada beda tingginya.
Melompat ke tempat yang lebih rendah memberi waktu di udara lebih lama,
sehingga jarak mendatar terjauhnya lebih jauh. Karena itu, pemeriksaan pada
bab ini menghitung jangkauannya untuk setiap pasangan tempat berdiri, bukan
memakai satu angka untuk seluruh level.

\subsection{Algoritmanya Sudah Anda Miliki}

Setelah jangkauan lompatan diketahui, pemeriksaannya menjadi persoalan yang
sudah Anda selesaikan pada Bab~X: pencarian jalur. Yang berubah hanya satu
hal.

\begin{kodekeluaran}
# Pada Bab X, tetangga sebuah petak adalah empat petak di sekelilingnya.
def tetangga(peta, dinding, petak):
    return [petak yang bersebelahan dan tidak berupa dinding]

# Pada Bab XVI, tetangga sebuah pijakan adalah pijakan lain yang dapat
# dicapai dengan berjalan, jatuh, atau melompat.
def tetangga(peta, dinding, daftar_pijakan, petak):
    hasil = [pijakan di kiri dan kanan yang dapat dicapai dengan berjalan]
    hasil += [pijakan lain yang dapat dicapai dengan melompat atau jatuh]
    return hasil
\end{kodekeluaran}

Sisanya sama persis: antrean prioritas, perkiraan jarak untuk mengarahkan
pencarian, dan tabel biaya. Itulah pelajaran yang perlu Anda bawa dari bab
ini: algoritma pencarian jalur bukan hanya untuk menggerakkan musuh, melainkan
juga untuk memeriksa level Anda sendiri. Yang menentukan kegunaannya adalah
aturan tetangganya.

\subsection{Tiga Pemeriksaan dan Hasilnya}

Skrip \code{ch16_uji_mandiri.py} menjalankan tiga pemeriksaan atas sebuah
level. Ketiga peta bawaan pada skrip itu memperlihatkan pemeriksaan yang lulus,
pemeriksaan yang menemukan tujuan yang tidak dapat dicapai, dan pemeriksaan
yang menemukan \istilah{soft-lock}.

\subsubsection*{Level yang bersih}

\begin{kodekeluaran}
$ python -I ch16_uji_mandiri.py --peta rapi
====================================================================
 UJI MANDIRI LEVEL  (tanpa pemain)
====================================================================
 Peta            : rapi  (30 x 12 petak)
 Tempat berdiri  : 44
 Tinggi lompat   : 4.25 petak
 Jarak mendatar  : 7.29 petak (pada ketinggian yang sama)

1. KETERJANGKAUAN
-----------------
 LULUS: seluruh tempat berdiri dapat mencapai tujuan.

2. JALAN KE ATAS
----------------
 Tanpa jalan ke atas: 25 dari 44 tempat berdiri.
   25 di antaranya setinggi atau lebih tinggi daripada
   tujuan, sehingga tidak perlu naik. Ini wajar.
 LULUS: tidak ada tempat di bawah tujuan yang hanya dapat turun.

3. SIMULASI TANPA PEMAIN
------------------------
 Titik awal diacak 200 kali.
 Berhasil mencapai tujuan: 200 (100.0%).
 Panjang jalur  : terpendek 1 pijakan, terpanjang 5 pijakan, rata-rata 2.6 pijakan.

KESIMPULAN
----------
 Level lolos seluruh pemeriksaan otomatis.
 Langkah berikutnya adalah playtest dengan pemain, bukan
 pemeriksaan otomatis. Yang dapat diperiksa mesin hanyalah
 apakah levelnya mungkin diselesaikan, bukan apakah levelnya
 menyenangkan.
\end{kodekeluaran}

\subsubsection*{Tujuan yang tidak dapat dicapai}

Pada peta kedua, tujuan diletakkan di atas tiang yang terlalu tinggi dan
terlalu jauh untuk dicapai. Program berjalan tanpa kesalahan apa pun; yang
salah hanyalah rancangan levelnya.

\begin{kodekeluaran}
$ python -I ch16_uji_mandiri.py --peta sempit
...
1. KETERJANGKAUAN
-----------------
 GAGAL: 44 dari 45 tempat berdiri (97.8%) tidak dapat mencapai tujuan.
...
2. JALAN KE ATAS
----------------
 PERIKSA: 21 tempat (46.7%) berada di bawah tujuan tetapi
 tidak dapat naik kembali.
...
 Jalan keluar terdekat: dari (10, 5) ke (20, 2).
 Diperlukan jarak 10.00 petak, sedangkan jangkauannya hanya 5.62 petak,
 yaitu kurang 4.38 petak. Perbaikan sekecil itu sudah cukup.

3. SIMULASI TANPA PEMAIN
------------------------
 Titik awal diacak 200 kali.
 Berhasil mencapai tujuan: 7 (3.5%).
\end{kodekeluaran}

Angka 3,5 persen itu penting untuk dipahami. Ia berarti permainannya masih
dapat diselesaikan dari sebagian kecil tempat awal, sehingga penguji yang
mencobanya beberapa kali mungkin menyimpulkan bahwa levelnya hanya
``sulit'', bukan bahwa levelnya tidak dapat diselesaikan.

\subsubsection*{\istilah{Soft-lock} yang tersembunyi}

Peta ketiga adalah yang paling berbahaya, sebab ia \emph{hampir} benar.
Tujuannya dapat dicapai dari hampir seluruh tempat, dan simulasi mengembalikan
angka keberhasilan yang tinggi. Yang salah hanyalah satu lubang:

\begin{kodekeluaran}
$ python -I ch16_uji_mandiri.py --peta terjebak
...
1. KETERJANGKAUAN
-----------------
 GAGAL: 5 dari 32 tempat berdiri (15.6%) tidak dapat mencapai tujuan.
 Ini calon soft-lock: pemain yang terlanjur sampai ke sana
 tidak dapat melanjutkan permainan, walaupun programnya tidak
 berhenti dan tidak ada pesan kesalahan apa pun.
   terjebak di petak (1, 10)
   terjebak di petak (2, 10)
   terjebak di petak (3, 10)
   terjebak di petak (4, 10)
   terjebak di petak (5, 10)
 Seluruh jalan keluar dari daerah terjebak terhalang
 dinding, sehingga tidak ada lompatan yang dapat diperiksa.
...
3. SIMULASI TANPA PEMAIN
------------------------
 Titik awal diacak 200 kali.
 Berhasil mencapai tujuan: 164 (82.0%).
\end{kodekeluaran}

Bandingkan kedua angka pada peta ketiga. Pemeriksaan keterjangkauan
mengembalikan 15,6 persen, sedangkan simulasi mengembalikan 82,0 persen.
Apabila Anda hanya menjalankan simulasinya, angka 82 persen akan terbaca
sebagai keberhasilan. Yang menyingkap \istilah{soft-lock}-nya adalah
pemeriksaan yang memeriksa \emph{setiap} tempat berdiri, bukan beberapa tempat
yang kebetulan terpilih.

\begin{catatabox}
Pelajarannya berlaku umum: pengujian yang bergantung pada kebetulan hanya
menemukan kesalahan yang kebetulan tersentuh. Pemeriksaan yang memeriksa
seluruh keadaan akan menemukan kesalahan yang jarang muncul, dan justru
kesalahan itulah yang paling merugikan pemain.
\end{catatabox}

\subsection{Apa yang Tidak Dapat Diperiksa Secara Otomatis}

Pemeriksaan otomatis pada bab ini memberi jawaban yang pasti atas satu
pertanyaan: apakah levelnya mungkin diselesaikan. Ia tidak dapat menjawab
pertanyaan lain, dan batas kemampuannya perlu Anda sadari sebelum Anda
mempercayainya.

\begin{itemize}
  \item Ia tidak dapat menilai apakah levelnya menyenangkan, adil, atau
        membosankan.
  \item Ia memakai jangkauan terbesar, yaitu pemain yang selalu berlari dengan
        kecepatan penuh dan melompat pada saat yang paling menguntungkan.
        Tempat yang lulus pemeriksaan ini belum tentu mudah dicapai pemain.
  \item Ia tidak memeriksa apakah bunyinya terdengar, apakah animasinya
        berganti, atau apakah antarmukanya menampilkan angka yang benar.
  \item Ia memeriksa level yang Anda tuliskan pada berkasnya, bukan level yang
        benar-benar Anda susun pada Unity. Penyesuaian tetap Anda lakukan
        sendiri.
\end{itemize}

Dua batas terakhir menjelaskan mengapa \istilah{playtest} tetap diperlukan,
dan itulah yang dibahas pada bagian berikutnya.
